\section{作品概述}
\label{sec:overview}

\subsection{功能与特性}
本作品在AMD Zynq-7100平台上集成摄像头采集、微小运动可视化、区域位移分析和交互控制。摄像头对准被测物后，系统连续处理图像，通过HDMI输出原画、运动放大画面或带频谱面板的处理画面，使观察者能够同时查看运动位置、运动形态和局部主频。

系统包含两个处理支路：整帧相位运动放大支路增强细微运动的可见性，固定感兴趣区域（ROI）测量支路从原始采集帧估计位移并计算时间频谱。用户通过串口屏选择视图、设置相位增益$\alpha$与曝光、控制启动和停止；板端依据寄存器读回和任务完成状态反馈执行结果。$\alpha$为0--50的无量纲相位调制参数，用于调节运动增强强度。

外部硬件采用开发板、相机和显示设备；设计重点为相位算法的PL实现、双支路帧调度、视频叠加和串口交互，形成从图像输入到结果显示的完整处理链路。

\subsection{应用领域}
本作品面向振动、信号处理和嵌入式系统教学。观察扬声器膜片等低频振动物体时，普通视频中的细微运动不易辨识，频谱曲线也难以独立说明运动发生的位置。系统以原画作为观察基准，以放大画面呈现运动形态，以ROI位移谱显示局部频率组成，可用于讲解周期振动、双频叠加、图像相位处理及软硬件协同。5\,Hz与10\,Hz扬声器实验构成了具体的演示场景。

工程应用方向为支架、薄板、悬臂梁及隔振结构的低频振动辅助观察。视觉采集无需在被测物上安装接触式传感器，板端处理与显示便于在实验现场查看结果。开展不同支撑条件的对照时，应固定相机视角、曝光、激励和成像比例，并结合参考传感器建立频率与物理位移的对应关系。

\subsection{主要技术特点}
本作品以相位运动处理方法\cite{wadhwa2013}为基础，将多尺度分解、局部相位处理、时间滤波和图像重建部署到PL，并设置独立的原始ROI位移分析支路，同时输出运动增强图像与局部频谱。

连续视频包含大量重复像素运算，并要求采集、计算和显示接口协调工作。设计利用HLS流水线、行缓存及RTL视频逻辑构建专用数据通路，由PS执行相机配置、帧调度和交互控制。DDR保存图像帧及部分历史状态，BRAM保存局部缓存与算法工作区；运算资源、存储容量和任务依赖按照实际数据访问方式组织。

主要算法采用浮点计算。双支路共享原始采集帧，分别维护完成状态；参数更新、结果提交和缓存复用均受相应完成条件约束，以保持处理数据与界面状态一致。总体关系见图\ref{fig:system}。

\subsection{主要性能指标}
表\ref{tab:overview-metrics}汇总系统配置、v18实板吞吐与2025.2实现资源，测试条件见第三部分。

\begin{table}[htbp]
  \centering\small
  \caption{主要参数与结果概览}
  \label{tab:overview-metrics}
  \begin{tabularx}{\linewidth}{@{}p{0.24\linewidth}Xp{0.17\linewidth}@{}}
    \toprule
    项目 & 参数或结果 & 结果来源 \\
    \midrule
    输入/显示 & $1280\times800$/1080p30 & 配置 \\
    ROI/时间窗 & $256\times256$/256点 & 配置 \\
    LUT/BRAM/DSP & 49.76\%/98.61\%/40.89\% & 2025.2基线 \\
    处理/显示速率 & 60.38/30\,frame/s & v18实板 \\
    \bottomrule
  \end{tabularx}
\end{table}

\subsection{主要创新点}
本作品围绕算法部署与系统集成形成以下工程创新：

（1）组合整帧运动放大与原始ROI位移谱，使运动形态和局部主频在同一终端对应呈现。

（2）以PL专用计算、片上工作区及DDR帧缓存承载算法，由PS协调双支路输入帧的使用和释放。

（3）将寄存器读回、处理完成和启停排空纳入交互，以设备执行状态驱动参数反馈。

\subsection{设计流程}
设计流程包括需求分析、算法参考、软硬件划分、HLS/RTL/PS实现、模块测试、硬件实现和板级联调，如图\ref{fig:development}所示。根据测试结果定位问题，返回相应模块修正后重新检查功能与资源。大模型辅助源码检索、调试、数据整理、报告与视频脚本编制，作者校正场景和需求；通过日志、数值核对及PDF检查确认结果，形成协作总结与工程指南。

\reportfig{fig1-1-development.pdf}{1.0}{系统研发与验证流程}{fig:development}
